\section{Experimental Setup}
\label{sec:experiments}

A central goal of this work is to show that competitive long-term forecasting is
achievable under a strict commodity-hardware budget. Accordingly, every
experiment in this paper---training, evaluation, ablation, and efficiency
profiling, for \method\ and all baselines---is run on a \emph{single} NVIDIA
RTX~3050~Ti laptop GPU with 4\,GB of VRAM. We re-implement and run all baselines
ourselves under identical data splits, normalization, seeds, and hardware, so
all comparisons are like-for-like rather than copied from other papers.

\subsection{Datasets}
\label{sec:exp:data}

We evaluate on seven widely used public long-term forecasting benchmarks: the
four Electricity Transformer Temperature datasets (ETTh1, ETTh2, ETTm1, ETTm2),
Weather, Electricity (ECL), and Traffic. The ETT datasets record seven
power-system variates; ETTh1/ETTh2 are sampled hourly and ETTm1/ETTm2 every
15 minutes. Weather records 21 meteorological variates at 10-minute resolution.
Electricity records 321 hourly electricity consumption variates; Traffic records
862 hourly road-occupancy variates from the Bay Area freeway system. The
higher-variate datasets (ECL, Traffic) expose the channel-mixing overhead of
Transformer-based models. Table~\ref{tab:datasets} summarizes all seven datasets.
All series are multivariate and forecast in the channel-independent setting
(Section~\ref{sec:method:setup}), except iTransformer and TimesNet which
are channel-mixing architectures (see Section~\ref{sec:exp:baselines}).
To probe behavior under distribution shift, we additionally use two datasets
outside the standard suite in the non-stationarity study
(Section~\ref{sec:results:nonstat}): Exchange-rate and ILI (national illness),
evaluated at the conventional shorter horizons $H \in \{24,36,48,60\}$.

\begin{table}[t]
\centering
\caption{Summary of the benchmark datasets. $C$ is the number of channels;
``Steps'' is the total series length. ETT splits follow the standard
12/4/4-month train/validation/test convention; all others use a chronological
70\%/10\%/20\% split.}
\label{tab:datasets}
\begin{tabular}{lcccc}
\toprule
Dataset     & $C$ & Frequency & Steps      & Split \\
\midrule
ETTh1       &  7  & hourly    & 17{,}420   & 12/4/4 mo \\
ETTh2       &  7  & hourly    & 17{,}420   & 12/4/4 mo \\
ETTm1       &  7  & 15-min    & 69{,}680   & 12/4/4 mo \\
ETTm2       &  7  & 15-min    & 69{,}680   & 12/4/4 mo \\
Weather     & 21  & 10-min    & 52{,}696   & 70/10/20\% \\
Electricity & 321 & hourly    & 26{,}304   & 70/10/20\% \\
Traffic     & 862 & hourly    & 17{,}544   & 70/10/20\% \\
\bottomrule
\end{tabular}
\end{table}

\subsection{Forecasting Protocol}
\label{sec:exp:protocol}

We follow the established long-term forecasting protocol
\citep{zeng2023dlinear,nie2023patchtst} so that our numbers are directly
comparable to published baselines. For the ETT datasets we use the canonical
12/4/4-month train/validation/test split; Weather uses a chronological
70\%/10\%/20\% split. Each series is z-score normalized using statistics
computed on the \emph{training portion only} (per channel), which prevents
leakage and matches the standard normalization that makes MSE/MAE comparable
across papers. The model's internal A-RevIN/RevIN then operates as an instance
normalization on top of this dataset-level normalization, and the loss and
metrics are computed in the train-z-scored space.

We report two lookback lengths: $L{=}336$, the strong setting for linear models,
as the \emph{headline} configuration, and $L{=}96$ for completeness. For every
lookback we forecast the four standard horizons $H \in \{96, 192, 336, 720\}$.
We evaluate with Mean Squared Error (MSE) and Mean Absolute Error (MAE),
averaged over all horizon steps and channels. Every $(\text{dataset}, H,
\text{model})$ cell is run with three seeds $\{2021, 2022, 2023\}$, and all
tables report the mean $\pm$ standard deviation over seeds.

\subsection{Baselines}
\label{sec:exp:baselines}

We compare against a sanity floor and seven competitive forecasters spanning the
linear, lightweight-frequency, and Transformer families:
\begin{itemize}
  \item \textbf{Naive (repeat-last)}: predicts the last observed value for the
        entire horizon; a parameter-free sanity floor.
  \item \textbf{NLinear} \citep{zeng2023dlinear}: subtracts the last lookback
        value, applies a single \texttt{Linear}$(L,H)$, and adds it back---a
        minimal distribution-shift-robust linear model.
  \item \textbf{DLinear} \citep{zeng2023dlinear}: moving-average trend/seasonal
        decomposition with two \texttt{Linear}$(L,H)$ heads, summed; the
        time-domain decomposition baseline \method\ generalizes.
  \item \textbf{RLinear} \citep{li2023rlinear}: RevIN \citep{kim2022revin} plus a
        single \texttt{Linear}$(L,H)$; the normalization baseline \arevin\
        generalizes.
  \item \textbf{FITS} \citep{xu2024fits}: RevIN plus a low-pass cutoff in the
        rFFT domain and a single complex linear layer; the closest lightweight
        frequency-domain prior art, which \emph{discards} the high-frequency band
        \method\ retains.
  \item \textbf{PatchTST-small} \citep{nie2023patchtst}: a deliberately small,
        channel-independent patch Transformer (patch length 16, stride 8,
        $d_{\text{model}}{=}64$, 4 heads, 2 layers, dropout 0.2). Run only on
        the five low-$C$ datasets (ETT + Weather); the attention memory of
        $O(L^{2})$ per patch makes it 4\,GB-infeasible on Electricity ($C{=}321$)
        and Traffic ($C{=}862$) at the lookbacks we use.
  \item \textbf{iTransformer} \citep{liu2024itransformer}: inverted self-attention
        over variate tokens (ICLR 2024); each channel is embedded as a token and
        attention operates across channels rather than time steps. Config:
        $d_{\text{model}}{=}128$, 4 heads, 2 encoder layers. A 4\,GB-feasible
        configuration for all seven datasets.
  \item \textbf{TimesNet-small} \citep{wu2023timesnet}: rFFT period detection plus
        2D Inception convolutions within each dominant period (ICLR 2023), scaled
        to $d_{\text{model}}{=}d_{\text{ff}}{=}32$, $k{=}3$ dominant periods,
        2 layers. As a channel-mixing model its memory scales moderately with $C$.
\end{itemize}
iTransformer and TimesNet are channel-mixing architectures; we run them on all
seven datasets with the same protocol but dataset-specific batch sizes (16 for
Weather and Electricity, 8 for Traffic) to remain within the 4\,GB budget.

\subsection{Implementation Details}
\label{sec:exp:impl}

All models are implemented in PyTorch 2.6.0 (CUDA 12.4) and trained on a single
RTX~3050~Ti laptop GPU (4\,GB). We use the Adam optimizer \citep{kingma2015adam}
with learning rate $10^{-3}$, no weight decay, and MSE loss. Training uses a
batch size of 32, at most 20 epochs with early stopping (patience 3 on the
validation loss), a type-1 step learning-rate schedule (halving the rate each
epoch after the first), and global gradient-norm clipping at 1.0. We train in
full \texttt{fp32} precision: the models are tiny, fp32 fits comfortably within
4\,GB, and it avoids the reproducibility noise of mixed precision. For
reproducibility we fix the \texttt{torch}, \texttt{numpy}, and \texttt{random}
seeds, enable deterministic algorithms, and disable cuDNN autotuning. At each
run we record parameters, analytic FLOPs, wall-clock training time per epoch, and
peak GPU memory via \texttt{torch.cuda.max\_memory\_allocated}; these populate the
efficiency analysis in Section~\ref{sec:efficiency}. Every reported number is
regenerable by a script in the repository.
